\section{Discussion}

The results support a more careful view of neural-symbolic closure discovery than the usual ``neural beats symbolic'' narrative. In matched-library settings, weak polynomial baselines are genuinely strong. This is not an embarrassing edge case for the neural model; it is exactly what one should expect when the true closure lies inside the baseline hypothesis class and the weak-form design matrix is well-conditioned.

The mismatch regime is where the neural stage becomes scientifically meaningful. Even there, however, the role of the neural model is not to guarantee exact law discovery. Its role is to provide a flexible numerical constitutive surrogate when the restricted prior library is misspecified. The surrogate is therefore an intermediate object: useful, but not yet a validated scientific law.

This is why the symbolic stage must be interpreted carefully. The experiments show that symbolic compression is highly effective as a distillation step. It can reduce constitutive curves to compact symbolic forms with very small additional surrogate error and almost no extra rollout degradation. But once that compression error is small, it ceases to be the dominant source of error. The remaining constitutive bias is dominated by mismatch and identification error from Stage~1. In this sense, the symbolic stage is not the main bottleneck; it is the place where the surrogate becomes interpretable, not the place where it becomes correct.

The noisy benchmarks sharpen this point rather than weaken it. Even with a weak-form training signal, Stage~1 closure recovery deteriorates markedly under $5\%$ observation noise, especially for the reaction term. That is not evidence against the paper's thesis. It is evidence for it: weak-form training mitigates spatial differentiation brittleness, but it does not eliminate the inverse-problem difficulty of constitutive recovery under noisy, weakly informative observations. The symbolic stage then faithfully inherits that degraded surrogate.

These observations suggest a clear agenda for future work. The main difficulty is reliable numerical constitutive recovery under limited excitation, sparse observations, and misspecified prior structure. Better Stage~1 recovery may come from richer excitation design, stronger identifiability diagnostics, adaptive trajectory selection, improved regularization, or more problem-aware surrogate classes. In a few matched-library sparse settings, the diffusion error improves slightly relative to the clean reference, which is consistent with a mild implicit smoothing effect rather than a genuine information gain. Once the numerical recovery problem is better controlled, the symbolic stage can remain simple and restricted.

The symbolic library is also intentionally narrow. It is designed to answer a specific question: can a learned surrogate be compressed into a compact family without changing its constitutive or dynamical behavior? It is not intended as a complete search space for nonlocal, singular, or history-dependent closures. If broader closure classes are needed, they should be introduced as controlled extensions of the same compress-and-reinsert protocol rather than as unvalidated formula mining.

The current study is intentionally narrow. It focuses on 1D periodic reaction-diffusion systems, synthetic data, and a controlled set of closure families. This keeps the mechanisms visible. Extending the same logic to higher-dimensional systems, transport-diffusion-reaction couplings, or more realistic observation models will introduce additional identifiability challenges, especially for anisotropic or multivariate closures. That broader program is important future work, but it should be built on the present lesson: low residual is not enough, and symbolic compression should be judged by what it preserves.
